\documentclass[10pt]{extarticle}

% ---------- Packages ----------
\usepackage{fancyhdr}
\usepackage{geometry}
\usepackage{hyperref}
\usepackage{xcolor}
\usepackage{amsmath, amsfonts, bm}
\usepackage{subcaption}
\usepackage{graphicx}
\usepackage{tikz}
\usepackage{nomencl}
\usepackage{array}
\usepackage{pdflscape}
\usepackage[export]{adjustbox}
\usepackage[utf8]{inputenc}
\usepackage{parskip}

\usetikzlibrary{shapes.geometric}
\usetikzlibrary{arrows.meta, arrows}

% ---------- Page setup ----------
\geometry{
    a4paper,
    total={170mm,257mm},
    left=20mm,
    top=20mm
}

% ---------- Headers and footers ----------
\pagestyle{fancy}
\pagenumbering{arabic}

\fancyhead[L]{Rocket Physics by Yaroslav Knyazkov}
\fancyfoot[C]{}
\fancyfoot[R]{\thepage}

% ---------- Hyperlinks ----------
\hypersetup{
    colorlinks=true,
    linkcolor={blue!50!black},
    citecolor={blue!50!black},
    urlcolor={blue!80!black}
}

% ---------- Equation numbering ----------
\numberwithin{equation}{subsection}

% ---------- Custom commands ----------
\newcommand{\todo}[1]{%
    \textcolor{red}{(#1)}%
}

\begin{document}

\tableofcontents

\newpage
\section{Current assumptions}\label{sec:assumptions}
This section collects, in one place, every simplification the current derivation and the
simulator rely on. Each is expected to be relaxed in a later phase; the relevant model or backlog
item is named where one exists, so a reader can tell ``current shortcut'' from ``permanent
choice'' at a glance.
\begin{itemize}
    \item \textbf{Planar (2D) flight, no crosswind.} Motion is confined to the vertical plane
    fixed by the launch azimuth $\chi$; the out-of-plane wind component is computed and discarded
    (\autoref{sec:wind}). No vertical wind either: wind velocity is horizontal only,
    updrafts/downdrafts are not modelled and are deliberately left for later, pending research
    into which model is actually worth the added data dependency (\autoref{sec:backlog}).
    \item \textbf{Gravity direction, parallel by default, with two non-parallel models
    available.} $\hat{z}_g$ can rotate with downrange distance, either treating the Earth as a
    sphere (radial) or as the WGS84 ellipsoid via its osculating sphere in the flight azimuth
    (elliptical) --- both in \autoref{sec:gravity} > Direction. Parallel, $\hat{z}_g\equiv
    \hat{z}_0$, is still the default, so every result produced with it is reproduced unchanged.
    Once a non-parallel model is enabled, the ground itself is curved
    (\autoref{eq:surface_point}) and the true altitude/impact condition follow
    \autoref{eq:h_true} rather than $h=-z$, $z=0$.
    \item \textbf{Earth rotation carried only in part, and in the fixed ground frame.} The
    \emph{static} effect is present --- the \emph{magnitude} of $g$ includes the centrifugal
    contribution through the WGS84 model (\autoref{sec:gravity}) --- and so is the in-plane
    Coriolis term in the translational equations of motion (\autoref{sec:noninertial}). What is
    dropped is the out-of-plane Coriolis component, which this 2D model has no degree of freedom
    to receive, and the residual centrifugal term, derived there but left unimplemented as some
    three orders of magnitude smaller. Separately, the decomposition of $\vec{\Omega}$ is frozen
    at the launch site's latitude and heading in $\mathcal{F}_0$ and does \emph{not} follow a
    rotating $\hat{z}_g$; the two corrections stay deliberately independent
    (\autoref{sec:noninertial}, closing caveat).
    \item \textbf{Constant moment arms (not constant mass).} Mass, centre of gravity and moment of
    inertia are \emph{not} held fixed: $m(t)$, $l_B(t)$ and $I(t)$ come from the propellant model
    of \autoref{sec:mass} and are carried by \autoref{sec:eom}. What is still held fixed is one
    step further out --- how the other two application points travel relative to a moving centre
    of gravity. The CoG-to-$D$ arm $l_1=l_D-l_B$ responds to $(\alpha,M)$ through the coefficient
    table (\autoref{sec:aero}) but not to CoG travel, and the CoG-to-$C$ thrust arm is a constant.
    \item \textbf{Thrust axis: fixed cant, no commanded control.} The nozzle cant $\delta$
    (\autoref{eq:thrust_frame}) and the lateral offsets $u_B,u_C$ (\autoref{sec:geometry}) are all
    carried as constant inputs --- as-built manufacturing quantities to be swept, not functions of
    the flight state. $\delta=u_B=u_C=0$ is the default and reproduces every earlier result
    unchanged. A \emph{commanded} deflection --- a gimbal and the control law driving it, the
    only pitch authority a rocket has off the rail --- is deliberately not modelled yet
    (\autoref{sec:backlog}).
    \item \textbf{No launch rail constraint, no recovery/parachute.} Free flight from $t=0$; see
    \autoref{sec:backlog}.
\end{itemize}

\newpage
\section{Rocket 2D}\label{sec:rocket2d}

\subsection{Frames}\label{sec:frames}
\begin{itemize}
    \item $\mathcal{F}_0$ --- ground/launch frame, origin at the launch point $O^*$.
    \item $\mathcal{F}_g$ --- gravity frame.
    \item $\mathcal{F}_r$ --- rocket (body) frame.
    \item $\mathcal{F}_t$ --- thrust frame.
\end{itemize}
$\mathcal{F}_0$ is fixed at the launch point $O^*$, with $\hat{x}_0$ horizontal (downrange) and
$\hat{z}_0$ down. The state $(x,z,\theta)$ is always integrated in this frame, regardless of which
gravity or ground model is switched on --- it is the one fixed thing in the whole document.
$\mathcal{F}_g$ carries the direction gravity actually points, $\hat{z}_g(\varphi)$, a function of
the downrange arc angle $\varphi$ introduced in \autoref{sec:gravity}. It coincides with
$\mathcal{F}_0$ by default, $\hat{x}_g(0)\equiv\hat{x}_0$ and $\hat{z}_g(0)\equiv\hat{z}_0$, and
only rotates away from it once a non-parallel direction model is enabled --- so $\mathcal{F}_g$
and $\mathcal{F}_0$ are kept notationally separate precisely because ``ground'' and ``the
direction gravity points'' are not always the same thing, even though they agree at the launch
point.

$\mathcal{F}_r$ is right-handed, with $\hat{x}_r$ along the rocket's nose axis (positive forward),
$\hat{z}_r$ completing the pitch plane with it, and $\hat{y}_r$ out of the pitch plane. Pitch
$\theta$ is measured from $\hat{x}_0$ (the ground-frame horizontal), positive nose-up, and the
pitch rate is $\omega=\dot\theta$, the rotation rate about $\hat{y}_r$. In ground-frame
components,
\begin{equation}\label{eq:body_axes_in_ground}
    \hat{x}_r = \cos\theta\,\hat{x}_0 - \sin\theta\,\hat{z}_0, \qquad
    \hat{z}_r = \sin\theta\,\hat{x}_0 + \cos\theta\,\hat{z}_0.
\end{equation}

$\mathcal{F}_t$ carries the actual line of thrust action, $\hat{x}_t$, which need not coincide
with the nose axis: a real nozzle can be canted, or otherwise misaligned, by a small angle
$\delta$ from the body axis --- a fixed property of how the motor is mounted, not of the flight
state. Giving thrust its own frame, rather than folding $\delta$ into every thrust term by hand,
is what lets that distinction be written down once:
\begin{equation}\label{eq:thrust_frame}
    \hat{x}_t = \cos\delta\,\hat{x}_r + \sin\delta\,\hat{z}_r, \qquad
    \hat{z}_t = -\sin\delta\,\hat{x}_r + \cos\delta\,\hat{z}_r, \qquad \hat{y}_t\equiv\hat{y}_r.
\end{equation}
For now $\delta$ is carried as a constant input --- a fixed manufacturing-type misalignment, swept
by hand to see how sensitive the airframe is to it. A commanded, time-varying gimbal angle driven
by a control law is future work (\autoref{sec:backlog}).

\subsection{Forces and moments}\label{sec:forces}
Weight is the main force the rocket has to overcome to increase its altitude:
\begin{equation}
    \vec{W} = mg\hat{z}_g,
\end{equation}
where $m$ is the rocket's mass and $g$ the gravitational acceleration along the local vertical
$\hat{z}_g$. Both factors are models rather than constants, and are chosen independently: the
magnitude $g$ and the direction $\hat{z}_g$ answer different questions --- how strong, and which
way --- and are supplied separately by \autoref{sec:gravity}.

Thrust is generated by ejecting propellant at high velocity, plus a contribution from the pressure
difference across the nozzle exit. With the thrust frame of \autoref{eq:thrust_frame} it is simply
\begin{equation}
    \vec{T} = T\hat{x}_t,
\end{equation}
purely axial in its own frame, with $T$ the thrust magnitude (\autoref{sec:thrust}); the cant
$\delta$ inside $\hat{x}_t$ carries the whole of the misalignment, so this equation never needs to
be rewritten as $\delta$ changes.

Motion through the surrounding air produces an aerodynamic force,
\begin{equation}
    \vec{A}_f = -A_x\hat{x}_r - A_z\hat{z}_r,
\end{equation}
with $A_x,A_z$ the drag and lift magnitudes, defined in the body frame (\autoref{sec:aero}). By
defining a reference point $D$ on the rocket, a pitching moment can also be written down:
\begin{equation}
    \vec{A}_m = -A_m\hat{y}_r,
\end{equation}
with $A_m$ the pitching-moment magnitude, about $\hat{y}_r$.

\subsection{Geometry}\label{sec:geometry}
Each of the three forces acts at its own point on the rocket. Every dimension below is defined in
the rocket frame $\mathcal{F}_r$ and measured from the nose tip $A$.

Weight is defined at the centre of gravity, here called $B$:
\begin{equation}
    \vec{AB} = -l_B\hat{x}_r - u_B\hat{z}_r.
\end{equation}
Thrust is defined at $C$, the point where the (possibly canted) thrust line of $\mathcal{F}_t$
crosses the body:
\begin{equation}
    \vec{AC} = -l_C\hat{x}_r - u_C\hat{z}_r.
\end{equation}
The aerodynamic force is defined at the reference point $D$:
\begin{equation}
    \vec{AD} = -l_D\hat{x}_r.
\end{equation}
The added $u$ dimensions refer to the axis misalignment: in a real rocket, the centre of gravity
and the nozzle can't be strictly on the axis, whether from manufacturing tolerance or from an
as-built mass asymmetry. As $D$ is a reference point rather than a measured one, no misalignment
is defined for it ($u_D\equiv0$); any asymmetry in the external surfaces instead shows up in the
aerodynamic coefficients themselves, e.g.\ in polar plots, not in a lateral offset of $D$. $u_B$
and $u_C$ are carried as ordinary constants alongside $l_B,l_C$ throughout what follows, on the
same footing as the thrust cant $\delta$ --- setting them to zero recovers the on-axis geometry
exactly, and there is no separate ``model'' for them beyond that: like $\delta$, they are inputs
representing an as-built offset, in principle obtained from manufacturing metrology of the actual
airframe rather than derived from anything else in this document.

\subsection{Equations of motion}\label{sec:eom}
\begin{gather}
    m\vec{\ddot{x}} = \sum \vec{F} = \vec{W} + \vec{T} + \vec{A}_f \label{eq:eom_trans}\\
    I\ddot{\theta} = \sum \vec{M}(B) \cdot \hat{y}_r = \left[\vec{BC}\times\vec{T} + \vec{BD}\times\vec{A}_f + \vec{A_m}\right]\cdot \hat{y}_r \label{eq:eom_rot}
\end{gather}
Moments are taken about the centre of gravity $B$, so weight contributes none:
$\vec{BB}\times\vec{W}=0$, for whatever direction $\vec{W}$ points. $m$ and $I$ are not constants
once the propellant model of \autoref{sec:mass} is active --- they become $m(t), I(t)$, and the
centre of gravity $B$ itself travels as $l_B(t)$. Reinterpreting $\mathcal{F}_0$ as Earth-fixed
rather than inertial adds one further translational term to \autoref{eq:eom_trans} and nothing at
all to \autoref{eq:eom_rot} (\autoref{sec:noninertial}). \autoref{sec:analysis} expands these two
equations by hand.


\newpage
\section{Analytical analysis}\label{sec:analysis}
This section works through the full configuration space of the pitching moment --- thrust on or
off, thrust axis aligned or misaligned --- and the static and dynamic stability results that hang
off it. Here ${}^r(\cdot)$ denotes a vector's components resolved along the body axes $\hat{x}_r,
\hat{z}_r$ at the current instant --- i.e.\ $m\,{}^r\ddot{x}=(m\ddot{\vec{x}})\cdot\hat{x}_r$ ---
\emph{not} a derivative taken within the rotating frame $\mathcal{F}_r$ itself (which would
additionally need that frame's own rotation-coupling terms). The equations below are therefore
just $\sum\vec{F}\cdot\hat{x}_r$, $\sum\vec{F}\cdot\hat{z}_r$ of \autoref{eq:eom_trans}, expanded,
and \autoref{eq:eom_rot} with the geometry of \autoref{sec:geometry} substituted in.

Two model-level quantities appear in the equations below and are worth naming up front, because
both are zero by default and both reduce this section to its simplest form when they are: the
thrust cant angle $\delta$ (\autoref{sec:forces}, \autoref{sec:thrust}), which is zero for an
ideally aligned nozzle, and the gravity-direction rotation angle $\varphi=x/R_e$
(\autoref{sec:gravity} > Direction), which is zero at the launch point and stays zero for the
whole flight if the local vertical is frozen there. Setting $\delta=\varphi=0$ recovers every
equation of this section in its simplest form.

\subsection{Forces}\label{sec:analysis_forces}
The aerodynamic force is a body-axis quantity already; thrust is one too once its own axis is
fixed relative to the body via $\mathcal{F}_t$ (\autoref{sec:forces}), and only weight has to be
resolved. Recall the body axes in ground-frame components, \autoref{eq:body_axes_in_ground}, with
$\theta$ measured from $\hat{x}_0$, positive nose-up --- the frame the simulator integrates $x,z$
in (\autoref{sec:frames}). Weight acts along the
\emph{local} vertical $\hat{z}_g(\varphi)=-\sin\varphi\,\hat{x}_0+\cos\varphi\,\hat{z}_0$, so
\begin{equation}
    \vec{W}\cdot\hat{x}_r = -mg\sin(\theta+\varphi), \qquad
    \vec{W}\cdot\hat{z}_r = +mg\cos(\theta+\varphi),
\end{equation}
i.e.\ the downrange rotation of gravity enters the body-axis force balance as nothing more than a
shift of the pitch angle, $\theta\to\theta+\varphi$. That is the whole of its effect, and it is
worth reading once: $\theta+\varphi$ is exactly the rocket's pitch above the \emph{local} horizon,
which dips below the launch-site horizontal by $\varphi$ at downrange distance $x=R_e\varphi$.
Adding thrust and the aerodynamic force,
\begin{align}
    m{}^r\ddot{x} &= -mg\sin(\theta+\varphi) + T\cos\delta - A_x \label{eq:analytic_force_x}\\
    m{}^r\ddot{z} &= mg\cos(\theta+\varphi) + T\sin\delta - A_z \label{eq:analytic_force_z}
\end{align}
The two limiting attitudes fix the signs, taking $\varphi=\delta=0$ for the check. Nose straight
up ($\theta=90^\circ$): weight is purely axial and purely retarding, $m{}^r\ddot{x}=-mg+T-A_x$,
with no normal component at all. Nose horizontal ($\theta=0$): $\hat{z}_r\equiv\hat{z}_g$, weight
is purely normal, $m{}^r\ddot{z}=mg-A_z$, and does not oppose the thrust in the slightest.

A third check fixes the sign of $\varphi$. Flying level in the launch frame ($\theta=0$) at
downrange $x>0$, \autoref{eq:analytic_force_x} gives $m{}^r\ddot{x}=-mg\sin\varphi<0$: gravity has
acquired a component that \emph{retards} the downrange motion. It should have --- the Earth's
centre sits below the launch point, so from a point downrange of it the pull is tilted backwards,
against the flight direction. At $x=10\,\mathrm{km}$ this is $g\varphi\approx
1.5\times10^{-2}\,\mathrm{m/s^2}$, and it grows linearly with downrange distance.

\subsection{Moment}\label{sec:analysis_moment}
Substituting $\vec{BC}=\vec{AC}-\vec{AB}$ and $\vec{BD}=\vec{AD}-\vec{AB}$ from
\autoref{sec:geometry} into \autoref{eq:eom_rot}, with $u_D\equiv0$,
$\hat{x}_r\times\hat{z}_r=-\hat{y}_r$, and the general thrust vector $\vec{T}=T\cos\delta\,
\hat{x}_r+T\sin\delta\,\hat{z}_r$ of \autoref{eq:thrust_frame}:
\begin{equation} \label{eq:moment_full}
    I\ddot{\theta} = (u_B - u_C)T\cos\delta + (l_C - l_B)T\sin\delta
                     - u_BA_x + (l_B - l_D)A_z - A_m
\end{equation}
Each term reads off the geometry directly. Thrust acts at $C$, whose displacement from the CoG is
$\vec{BC}=(l_B-l_C)\hat{x}_r+(u_B-u_C)\hat{z}_r$: its lateral offset $u_B-u_C$ gives an arm to the
\emph{axial} thrust component $T\cos\delta$, and its axial offset $l_C-l_B$ gives an arm to the
\emph{normal} thrust component $T\sin\delta$. The aerodynamic force acts at $D$, so the axial arm
$l_B-l_D$ carries its normal component $A_z$ and the lateral arm $u_B$ carries its axial component
$A_x$. Weight does not appear here at all, and cannot: moments are taken about the centre of
gravity $B$, so $\vec{BB}\times\vec{W}=0$ for \emph{every} orientation of $\vec{W}$ --- rotating
the local vertical (\autoref{sec:gravity} > Direction) is therefore a purely translational
correction, exactly like the Coriolis term of \autoref{sec:noninertial}, and for the same reason
(a body force distributed in proportion to mass has its resultant at the CoG).

$l_C$ is a live quantity here in a way it is not for on-axis thrust. An on-axis thrust
$\vec{T}=T\hat{x}_r$ applied at a point whose displacement from the CoG has axial part
$(l_B-l_C)\hat{x}_r$ contributes $(l_B-l_C)T\,\hat{x}_r\times\hat{x}_r=0$ from that part: sliding
the application point forward or aft along the axis changes nothing, because the force is parallel
to the displacement it would act through. Only the lateral offset $u_C$ survives in that case.
Once $\delta\neq0$ the thrust has a component normal to $\hat{x}_r$, and a normal force
\emph{does} care how far along the axis it is applied --- hence the second term of
\autoref{eq:moment_full}, and hence the axial position of the nozzle mattering for the moment.

\autoref{eq:moment_full} specialises four ways, and the following subsections take them in turn.
Throughout, \emph{aligned} is shorthand for the on-axis geometry $u_B=u_C=0$ \emph{and}
$\delta=0$; $u_B,u_C,\delta$ are all carried as ordinary constant inputs (\autoref{sec:geometry},
\autoref{sec:thrust}), so setting all three to zero recovers this case exactly.
\begin{center}
\begin{tabular}{l l l}
 & \textbf{aligned} ($\delta=0$) & \textbf{misaligned} ($\delta\neq0$) \\
\textbf{no thrust} ($T=0$) & \autoref{eq:moment_nothrust} & identical --- $T=0$ scales $\delta$ away \\
\textbf{with thrust} ($T>0$) & \autoref{eq:moment_aligned} & \autoref{eq:moment_misaligned} \\
\end{tabular}
\end{center}
Only three of the four cells are distinct, and only one of them is new.

\subsubsection{No thrust: coast and descent}\label{sec:analysis_nothrust}
The coast and descent phases are worth isolating from the boost phase. Setting $T=0$ drops the
thrust term and leaves the airframe's purely aerodynamic behaviour, free of any dependence on the
motor's thrust level or on where the thrust acts --- so whatever stability it shows is a property
of the geometry and the aerodynamic coefficients alone, not of the motor:
\begin{equation} \label{eq:moment_nothrust}
    I\ddot{\theta} = -u_BA_x + (l_B - l_D)A_z - A_m.
\end{equation}
Thrust misalignment is physically irrelevant in this case rather than merely neglected: $T=0$
multiplies both $\cos\delta$ and $\sin\delta$ out of \autoref{eq:moment_full}, so the two
right-hand columns of the table above collapse onto one another exactly, for any $\delta$
whatever. It is stated rather than passed over silently because it is the reason the
static-stability result below --- which is built on this case --- needs no revision when the
misalignment model is switched on.

\subsubsection{With thrust, aligned}\label{sec:analysis_aligned}
If the thrust axis and the body axis are aligned ($u_B=u_C=0$, $\delta=0$), the equation turns
into
\begin{equation} \label{eq:moment_aligned}
    I\ddot{\theta} = (l_B - l_D)A_z - A_m.
\end{equation}
With the centre of gravity and the nozzle both on the body axis and the thrust along it, the first
two terms of \autoref{eq:moment_full} vanish together, the remaining lateral arm $u_B$ of
\autoref{eq:moment_nothrust} vanishes too, and the pitching behaviour reduces to the normal force
acting through the arm $l_B-l_D$ against the airframe's own pitching moment. The notable feature
of this case is what is \emph{missing}: the motor. Thrust has dropped out of the pitching equation
entirely, so an aligned rocket's attitude dynamics are the same during boost as during coast,
however large $T$ is --- not as an approximation, but as an exact consequence of the geometry it
is given.

\subsubsection{With thrust, misaligned}\label{sec:analysis_misaligned}
With the nozzle canted by a small fixed $\delta$ and the geometry otherwise on-axis
($u_B=u_C=0$):
\begin{equation} \label{eq:moment_misaligned}
    I\ddot{\theta} = \underbrace{(l_C - l_B)T\sin\delta}_{M_\delta} + (l_B - l_D)A_z - A_m.
\end{equation}
The new term $M_\delta$ is the point of modelling $\delta$ at all, and it behaves quite unlike
everything else in this section. Every other contribution to $\ddot{\theta}$ is aerodynamic and
therefore vanishes with $q$ and with $\alpha$; $M_\delta$ does not. It survives $A_z=A_m=0$, it
survives $\alpha=0$, and it survives $q=0$. For as long as the motor burns, the airframe carries a
built-in disturbance torque proportional to the thrust it is producing --- largest, in absolute
terms, exactly at the thrust peak, and largest \emph{relative} to the restoring moment just off
the pad and again high up, where $q$ is small.

Its sign is worth pinning down once. $l_C>l_B$ for any real rocket (the nozzle is aft of the
centre of gravity), and $\hat{z}_r$ points towards the rocket's belly (\autoref{sec:frames}), so a
positive $\delta$ pushes the tail towards $+\hat{z}_r$ and pitches the nose \emph{up}, $M_\delta>0$.

$\delta$ is a constant, so $\partial M_\delta/\partial\alpha=0$, and \autoref{eq:dtheta_dalpha} ---
the derivative the static margin of \autoref{eq:static_margin} is built on --- is untouched. A
misalignment adds a constant \emph{offset} to $\ddot{\theta}$, not a slope: it shifts the
equilibrium the airframe settles at, without changing whether that equilibrium is stable.
Concretely, setting $\ddot{\theta}=0$ in \autoref{eq:moment_misaligned} with the linearised
coefficients of \autoref{eq:aero_linearised} gives the trim angle of attack,
\begin{equation} \label{eq:alpha_trim}
    \alpha_{trim} = -\frac{M_\delta}{I\,\partial\ddot{\theta}/\partial\alpha}
        = \frac{-(l_C - l_B)T\sin\delta}
               {qS\bigl[(l_B - l_D)c_{z,\alpha} - c_{m,\alpha}d_{ref}\bigr]},
\end{equation}
taking $c_{z,0}=c_{m,0}=0$ for a body of revolution. A statically stable airframe has a negative
denominator (\autoref{eq:dtheta_dalpha}), so $\alpha_{trim}$ carries the sign of $M_\delta$: the
rocket does not diverge, it flies slightly nose-off the relative wind, at a constant small angle
whose normal force then curves the trajectory. That curvature, not a tumble, is the observable
consequence of a manufacturing-grade misalignment on a stable rocket.

Two limits are worth reading off \autoref{eq:alpha_trim} rather than discovering in a run.
$\alpha_{trim}\propto T/q$, so it blows up wherever there is thrust without dynamic pressure:
immediately after ignition, and during any high-altitude burn. And it blows up outright as the
static margin approaches zero, since the denominator is $I$ times \autoref{eq:dtheta_dalpha} --- a
marginally stable rocket has no authority to trim against even a tiny built-in torque. Neither
limit is a defect of the model; both are the physical reason launch rails exist.

The lateral offsets $u_B,u_C$ (\autoref{sec:geometry}) sit alongside $\delta$ in
\autoref{eq:moment_full} on exactly the same footing --- constant inputs, not yet exercised in
this section's worked cases, which all set $u_B=u_C=0$ to isolate $\delta$ on its own. Re-running
this section with $u_B,u_C\neq0$ is a direct substitution into \autoref{eq:moment_full}, adding a
further constant offset $(u_B-u_C)T\cos\delta-u_BA_x$ alongside $M_\delta$; it is not carried
through here only because there is no case yet where isolating it separately from $\delta$ is
useful. A \emph{commanded}, time-varying gimbal $\delta(t)$ driven by a control law remains future
work (\autoref{sec:backlog}).

\subsection{Static stability}\label{sec:analysis_stability}
Static stability is a steady measure of the rocket's stability against a disturbance in angle of
attack $\alpha$. Use the no-thrust, axis-aligned case --- \autoref{eq:moment_nothrust} with
$u_B=0$, which is the same right-hand side as \autoref{eq:moment_aligned}
(\autoref{sec:analysis_moment}: with $u_B=u_C=\delta=0$ the thrust contributes nothing, so the
boost and coast phases share one pitching-moment equation). Before the analysis, expand and
linearise $A_z,A_m$:
\begin{equation} \label{eq:aero_linearised}
    A_z = qc_zS,\quad A_m = qc_mSd_{ref},
\end{equation}
where $q$ is dynamic pressure, $c_z,c_m$ the lift and pitching-moment coefficients (written $C_N$
and $C_m$ in \autoref{sec:aero}; the lowercase form is kept here purely for historical reasons),
and $S,d_{ref}$ the reference area and length shared with \autoref{sec:aero}. Coefficients $c_z,
c_m$ are linearised as
\begin{align}
    c_z &= c_{z,0} + \alpha c_{z,\alpha}, \\
    c_m &= c_{m,0} + \alpha c_{m,\alpha}.
\end{align}
By taking the partial derivative of the angular acceleration with respect to $\alpha$, we expect
$\partial\ddot\theta/\partial\alpha<0$: a disturbance in $\alpha$ should produce an angular
acceleration that reduces it, i.e.\ the rocket stabilises the trajectory rather than diverging
from it. The expansion and derivative give
\begin{equation} \label{eq:dtheta_dalpha}
    \frac{\partial \ddot\theta}{\partial \alpha} = \frac{qS}{I}\biggl[(l_B - l_D)c_{z,\alpha} - c_{m,\alpha}d_{ref}\biggr],
\end{equation}
which is negative, provided $c_{z,\alpha}>0$ (the normal force grows with angle of attack --- the
normal case, and the only one in which dividing through by $c_{z,\alpha}$ preserves the sense of
the inequality), when the static margin satisfies
\begin{equation} \label{eq:static_margin}
    l_B - l_D < d_{ref}\frac{c_{m,\alpha}}{c_{z,\alpha}}.
\end{equation}
Both lengths are measured from the nose tip $A$ (\autoref{sec:geometry}), so $l_B-l_D$ is
\emph{negative} exactly when the aerodynamic reference point $D$ sits aft of the centre of gravity
$B$ --- the familiar ``centre of pressure behind centre of gravity'' rule. In the limiting case
where $D$ is the centre of pressure itself, $c_{m,\alpha}=0$ by definition and
\autoref{eq:static_margin} collapses to precisely that rule.

\paragraph{Static margin, in calibers.} $l_B-l_D$ is a dimensional length and, on its own, not
comparable between airframes of different size. The rocketry convention is instead to report the
static margin non-dimensionally, divided by the reference length $d_{ref}$ (the body diameter),
and with a sign flipped so that a \emph{stable} rocket gets a \emph{positive} number:
\begin{equation}
    \mathrm{SM} = \frac{l_D-l_B}{d_{ref}} = \frac{l_1}{d_{ref}},
\end{equation}
using $l_1=l_D-l_B$ exactly as already defined in \autoref{sec:aero}, so no new simulator quantity
is introduced --- only its non-dimensional report changes. $\mathrm{SM}>0$ now matches the
familiar caliber convention directly: the CoP sits $\mathrm{SM}$ body-diameters aft of the CoG for
a stable rocket, and \autoref{eq:static_margin} is equivalent to
$\mathrm{SM}>-c_{m,\alpha}/c_{z,\alpha}$.

\subsubsection{Dynamic stability}\label{sec:analysis_dynamic}
Everything above is a \emph{static} criterion: it asks only whether a disturbance in $\alpha$
produces a restoring moment, not what the resulting oscillation actually does next. That question
needs a fourth coefficient, the \textbf{pitch-damping derivative} $c_{m,q}$, the sensitivity of the
pitching-moment coefficient to a non-dimensional pitch rate $\hat{q}=\omega d_{ref}/(2V)$:
\begin{equation}
    c_m = c_{m,0} + \alpha c_{m,\alpha} + \hat{q}\,c_{m,q}.
\end{equation}
$c_{m,q}$ arises purely from the body's own rotation, exactly as its name suggests. When the
rocket rotates about the centre of gravity at rate $\omega$, every point on the body picks up an
extra velocity $\vec{\omega}\times\vec{r}_i$ relative to the CoG; a point a distance $l_i$ aft of
the CoG therefore sees a local angle of attack that differs from the CoG's own $\alpha$ by
$\Delta\alpha_i\approx+\omega l_i/V$. Fins mounted well aft of the CoG --- exactly where a stable
airframe puts its stabilising surfaces --- see the largest such shift, and integrating the
resulting change in the local normal-force distribution against the moment arm produces a moment
proportional to $\omega$: a true \emph{damping} term, present even when $\alpha$ itself is zero.
Note the sign convention here is the \emph{opposite} of the usual aerospace one: this document
fixed positive $A_m$ (and hence positive $c_m$) to pitch the nose \emph{down}
(\autoref{sec:forces}, $\vec{A}_m=-A_m\hat{y}_r$ feeding $I\ddot\theta=\dots-A_m$), which is why
\autoref{eq:dtheta_dalpha} already needs $c_{m,\alpha}>0$, not the textbook $C_{m\alpha}<0$, for
static stability. The same flip applies here: a conventional, aft-finned, statically stable
airframe has $c_{m,q}>0$ in this document's convention --- a positive (nose-up) pitch rate then
produces an \emph{additional} positive (nose-down) moment contribution, opposing the rotation that
caused it.

With this term, \autoref{eq:moment_aligned} linearises to
\begin{equation}
    I\ddot\theta = qS(l_B-l_D)c_{z,\alpha}\,\alpha - qSd_{ref}c_{m,\alpha}\,\alpha
        - \frac{qSd_{ref}^2}{2V}c_{m,q}\,\omega,
\end{equation}
(again taking $c_{z,0}=c_{m,0}=0$), which separates into a restoring (``stiffness'') term and a
damping term,
\begin{equation}
    M_\alpha \equiv \frac{\partial\ddot\theta}{\partial\alpha}
        = \frac{qS}{I}\bigl[(l_B-l_D)c_{z,\alpha}-d_{ref}c_{m,\alpha}\bigr], \qquad
    M_q \equiv \frac{qSd_{ref}^2}{2VI}c_{m,q},
\end{equation}
where $M_\alpha$ is exactly \autoref{eq:dtheta_dalpha}, negative for a statically stable airframe,
and $M_q$ is positive precisely when $c_{m,q}>0$ (the damping sign of this document's convention,
above).
Taking the short-period approximation --- the flight-path angle changes slowly enough over one
oscillation that $\dot\alpha\approx\dot\theta=\omega$ --- gives a standard damped-oscillator
equation for the angle-of-attack perturbation,
\begin{equation}
    \ddot\alpha + M_q\dot\alpha - M_\alpha\alpha = 0,
\end{equation}
with natural frequency and damping ratio
\begin{equation}
    \omega_n = \sqrt{-M_\alpha}, \qquad \zeta = \frac{M_q}{2\omega_n} = \frac{M_q}{2\sqrt{-M_\alpha}},
\end{equation}
valid only when $M_\alpha<0$ (statically stable); if it is not, the disturbance grows rather than
oscillates, and $\omega_n,\zeta$ are not meaningful. This is the framework the dynamic-stability
question needs; it is not yet quantitative, because $c_{m,q}$ is not supplied by
\autoref{sec:aero}'s coefficient database or its fallback. Getting it requires either a measured
value (forced-oscillation wind-tunnel testing, or an unsteady CFD run) or a full distributed-force
derivation from the airframe's own geometry along the lines sketched above; neither is done here.
\todo{needs a value or model for $c_{m,q}$ before this subsection is quantitative}


\newpage
\section{Models}\label{sec:models}
Each subsection below gives the interchangeable models behind one symbol of
\autoref{sec:rocket2d}. They are ordered by \emph{force}, following \autoref{sec:forces}: first
the three forces themselves --- weight (\autoref{sec:gravity}), thrust (\autoref{sec:thrust}),
together with the mass it consumes and the centre-of-gravity travel that follows
(\autoref{sec:mass}), and the aerodynamic force (\autoref{sec:aero}) --- and then the shared
environment the last of them is computed from: the air (\autoref{sec:atmosphere}), its motion
(\autoref{sec:wind}), the relative flow those two produce (\autoref{sec:relwind}), and finally the
fictitious-force corrections that come from the frame rather than from any force at all
(\autoref{sec:noninertial}). This ordering deliberately allows forward references: a model
sometimes uses a quantity defined further down (\autoref{sec:thrust}'s pressure correction needs
$p_{amb}$ from \autoref{sec:atmosphere}, and \autoref{sec:aero} needs $\rho$ and $q$ from
\autoref{sec:atmosphere} and $\alpha$ from \autoref{sec:relwind}); every such symbol is
cross-referenced at the point of use.

\subsection{Gravity}\label{sec:gravity}
Interchangeable models for the gravitational acceleration magnitude $g$ and direction $\hat{z}_g$
used in $\vec{W}=mg\hat{z}_g$ (\autoref{sec:forces}). Magnitude and direction are chosen
independently, since they answer different questions --- how strong, and which way. $h=-z$ is
altitude above the launch point, $R_e$ the Earth's mean radius.

\subsubsection{Magnitude}
\begin{itemize}
    \item \textbf{Constant} (Phase 1 baseline): $g=g_0$.
    \item \textbf{Radial} (inverse-square): exact for a spherically symmetric, non-rotating Earth,
    \begin{equation}
        g(h)=g_0\left(\frac{R_e}{R_e+h}\right)^2.
    \end{equation}
    Here $R_e+h$ is the distance from the Earth's centre, so $h$ is height above the
    \emph{local} spherical surface --- the same reading of the coordinates the direction model
    below uses.
    \item \textbf{Real data (WGS84)}: Somigliana normal gravity at the ellipsoid surface, plus the
    second-order free-air correction to height above it,
    \begin{gather}
        \gamma_0(\phi) = \gamma_e\frac{1+k\sin^2\phi}{\sqrt{1-e^2\sin^2\phi}} \\
        g(\phi, h) = \gamma_0\left[1 - \frac{2}{a_e}\left(1+f+m-2f\sin^2\phi\right)h + \frac{3}{a_e^2}h^2\right],
    \end{gather}
    with $\phi$ the launch site's geodetic latitude (fixed for the flight), $h$ the altitude above
    the ellipsoid, and published WGS84 constants $a_e=6378137\,\mathrm{m}$,
    $f=1/298.257223563$, $m=0.00344978650684$, $\gamma_e=9.7803253359\,\mathrm{m/s^2}$,
    $k=0.00193185265241$, $e^2=0.00669437999013$. This variant already includes the Earth's
    rotation (centrifugal) contribution to the \emph{magnitude} of $g$ through $\gamma_0(\phi)$ ---
    a separate, static effect from the dynamic (Coriolis) correction of
    \autoref{sec:noninertial}. A gridded dataset (e.g.\ EGM2008) would resolve anomalies two
    orders of magnitude below this model's own error and is not used here.
\end{itemize}

\subsubsection{Direction}
True gravitational acceleration points radially inward, toward the Earth's centre, so the local
vertical is not the same direction at two points a finite distance apart. Three variants, in
increasing fidelity:
\begin{itemize}
    \item \textbf{Parallel} (the default): $\hat{z}_g\equiv\hat{z}_0$, frozen at the launch point
    $O^*$'s own local vertical for the whole flight.
    \item \textbf{Radial} (spherical, non-parallel): the Earth is modelled as a sphere of radius
    $R_e$ centred at $O$ --- a point distinct from the launch point $O^*$, which merely lies on
    the sphere.
    \item \textbf{Elliptical} (non-parallel, origin still at $O$): the same WGS84 ellipsoid as the
    Magnitude model above, whose local vertical is the ellipsoid normal rather than the direction
    to $O$.
\end{itemize}

\paragraph{Construction (radial).} Work in the launch frame $\mathcal{F}_0$, origin at the launch
point $O^*$, axes $\hat{x}_0$ (downrange, horizontal) and $\hat{z}_0$ (down). The sphere's centre
$O$ then lies at $(0,R_e)$ in those coordinates --- directly below $O^*$ along $\hat{z}_0$, since
$\hat{z}_0$ is the down direction there. Parametrise the spherical surface by the arc angle
$\varphi$ subtended at $O$, measured from $O^*$ towards $+\hat{x}_0$:
\begin{equation} \label{eq:surface_point}
    \vec{O^*P}(\varphi) = R_e\sin\varphi\,\hat{x}_0 + R_e(1-\cos\varphi)\,\hat{z}_0,
\end{equation}
which is the launch point itself at $\varphi=0$ and drops away below the launch-point tangent
plane by $R_e(1-\cos\varphi)\approx x^2/2R_e$ thereafter --- the Earth's curvature, used below for
the ground and impact condition, though not, by default, for anything else. The local ``down'' at
$\vec{P}(\varphi)$ is the unit vector from the surface point to the centre, $(\vec{O}-\vec{P})/R_e$,
and the local downrange (horizontal) direction is the normalised surface tangent
$R_e^{-1}\,\mathrm{d}\vec{P}/\mathrm{d}\varphi$:
\begin{gather}
    \hat{x}_g(\varphi) = \cos\varphi\,\hat{x}_0 + \sin\varphi\,\hat{z}_0
        \label{eq:xhat_g_rotated} \\
    \hat{z}_g(\varphi) = -\sin\varphi\,\hat{x}_0 + \cos\varphi\,\hat{z}_0
        \label{eq:zhat_g_rotated}
\end{gather}
This is a rigid rotation of the launch-point basis about the axis normal to the flight plane: the
pair is orthonormal for every $\varphi$, has determinant $+1$ (so the triad stays right-handed and
$\hat{y}_g\equiv\hat{y}_0$ is untouched), and reduces to $\hat{x}_0,\hat{z}_0$ at $\varphi=0$. Note
which way it turns: the local horizon \emph{dips} (positive $\hat{z}_0$, i.e.\ downward) as one
moves downrange, and the local vertical tilts backwards, away from the flight direction. The angle
itself is taken as
\begin{equation} \label{eq:phi_downrange}
    \varphi = \frac{x}{R_e},
\end{equation}
with $x$ the simulation's own downrange coordinate --- reading $x$ as ground arc length, the same
reading the ground track $(x\sin\chi, x\cos\chi)$ of \autoref{sec:atmosphere} > Field already uses.

\paragraph{Weight.} The weight vector becomes, in launch-frame components,
\begin{equation} \label{eq:weight_rotated}
    \vec{W} = mg\,\hat{z}_g(\varphi)
        = mg\left(-\sin\varphi\,\hat{x}_0 + \cos\varphi\,\hat{z}_0\right),
\end{equation}
with $g$ from whichever magnitude model above is selected. That is the complete change: the
downrange channel gains $-mg\sin\varphi$, the vertical channel is scaled by $\cos\varphi$, and
nothing else in the equations of motion is touched. In particular the rotational equation
\autoref{eq:eom_rot} is untouched for a reason that holds exactly and not just to leading order
--- moments are taken about the centre of gravity, so weight contributes none whatever direction
it points (\autoref{sec:analysis_moment}). In body axes the whole effect collapses to the
substitution $\theta \to \theta + \varphi$ in the weight terms
(\autoref{eq:analytic_force_x}, \autoref{eq:analytic_force_z}).

\paragraph{Construction (elliptical).} Unlike the sphere, the ellipsoid's surface normal at
geodetic latitude $\phi$ does \emph{not} pass through the centre $O$ except along the equator and
at the poles: $O$ here anchors the coordinate system that defines $\phi$ (the same $\phi$ the
Magnitude model uses), it is not literally the point the local vertical points towards. To reuse
the rotation construction above rather than build a separate one, the ellipsoid is approximated
near the launch point $O^*$ by its \emph{osculating sphere} in the flight-plane azimuth $\chi$:
the sphere tangent to the ellipsoid at $O^*$ whose curvature matches the ellipsoid's own curvature
in that particular direction. Its radius is given by Euler's (Meusnier's) formula for the normal
curvature of a surface of revolution,
\begin{equation}
    R_\chi(\phi_0,\chi) = \left(\frac{\cos^2\chi}{M(\phi_0)} + \frac{\sin^2\chi}{N(\phi_0)}\right)^{-1},
\end{equation}
with $M,N$ the meridional and prime-vertical radii of curvature of the same WGS84 ellipsoid
($a_e,e^2$ as above),
\begin{equation}
    M(\phi) = \frac{a_e(1-e^2)}{(1-e^2\sin^2\phi)^{3/2}}, \qquad
    N(\phi) = \frac{a_e}{\sqrt{1-e^2\sin^2\phi}},
\end{equation}
and $\phi_0,\chi$ the launch site's geodetic latitude and heading. Every formula of the radial
construction above then carries over unchanged with $R_e\to R_\chi(\phi_0,\chi)$: $\varphi =
x/R_\chi$, and $\hat{x}_g(\varphi),\hat{z}_g(\varphi)$ exactly as in
\autoref{eq:xhat_g_rotated}--\autoref{eq:zhat_g_rotated}. This matches the true ellipsoid normal to
the same order the osculating-sphere approximation itself is good to: exact at $O^*$ and correct
to second order in $\varphi$ nearby, departing from the true normal at third order and beyond,
because the ellipsoid's own curvature keeps changing as the geodetic latitude drifts with
downrange distance --- a refinement not carried here.

\paragraph{Ground and impact condition.} Both non-parallel variants make the ground itself
curved, \autoref{eq:surface_point}, so the flat-Earth altitude $h=-z$ used everywhere else in this
document is only an approximation once one of them is switched on. The rocket's own position is
$\vec{O^*R}=x\hat{x}_0+z\hat{z}_0$; projecting its offset from the surface point at the
\emph{same} downrange arc angle onto the local ``up'' direction $-\hat{z}_g(\varphi)$ gives the
true altitude above the curved ground,
\begin{equation} \label{eq:h_true}
    h_{true}(x,z) = \left[\vec{O^*R}-\vec{O^*P}(\varphi)\right]\cdot\bigl(-\hat{z}_g(\varphi)\bigr)
        = x\sin\varphi - z\cos\varphi - R_e(1-\cos\varphi), \qquad \varphi=\frac{x}{R_e}.
\end{equation}
For small $\varphi$ this reduces to $h_{true} = x^2/2R_e - z + O(x^4/R_e^3)$: larger than the flat
reading $-z$ by exactly the curvature drop quoted above (the ground has dropped away beneath a
rocket that has not, so the true clearance is larger). \autoref{eq:h_true} is the altitude to use,
in place of $-z$, for every altitude-dependent model once a non-parallel direction model is
active --- the atmosphere and wind tiers, and the radial/WGS84 magnitude models --- and the impact
condition becomes $h_{true}(x,z)=0$ in place of $z=0$. For the elliptical variant the same
substitution $R_e\to R_\chi(\phi_0,\chi)$ applies throughout \autoref{eq:h_true} as well.

\subsubsection{Vector field (magnitude + direction)}
Data from some dataset, using the same station/field interpolation as \autoref{sec:atmos_tiers}.

\subsection{Thrust}\label{sec:thrust}
\subsubsection{Magnitude}
Magnitude of thrust is the result of conservation of momentum (ejection of fuel) and pressure
differences, which can be computed analytically for an internal-combustion nozzle:
\begin{equation}
    T = \dot{m}v_e + A_e(p_e - p_{atm}).
\end{equation}
In real life, however, the data from a static test is used instead. It is given as a tabulated
curve, linearly interpolated between measured $(t_i,T_i)$ points, and held at $T_1$ for $t<t_1$
and zero for $t>t_N$; burnout is the falling zero-crossing of this curve:
\begin{equation}
    T_{ref}(t) = \mathrm{lerp}\!\left(t; \{t_i\}, \{T_i\}\right).
\end{equation}
As the test is static, there is no variation in atmospheric pressure. To account for the change in
altitude, the net thrust once the motor is flowing is
\begin{equation}
    T(t, h) = T_{ref}(t) + A_e\left(p_{ref} - p_{amb}(h)\right),
\end{equation}
valid only while $T_{ref}(t)>0$ (i.e.\ while the motor is actually flowing); $p_{amb}(h)$ comes
from \autoref{sec:atmosphere}, $A_e$ is the nozzle exit area, and $p_{ref}$ is the ambient pressure
the curve was measured at (the test-stand pressure, not necessarily the launch site's own pad
pressure). $A_e=0$ recovers $T_{ref}(t)$ exactly.

\subsubsection{Direction}
Done, in this document, by a fixed misalignment rather than by active thrust control. The nozzle
cant $\delta$ of \autoref{eq:thrust_frame} is carried as a constant input --- a manufacturing-type
error, chosen once per airframe, not a function of $t$ or $h$ --- and is zero by default. The two
things this document keeps deliberately separate are: $\delta$ does not scale any part of the
magnitude model above, and the mass flow of \autoref{sec:mass} is driven by $T_{ref}(t)$
regardless of which way the nozzle points. The only coupling is the obvious one, that the axial
component actually delivered is $T\cos\delta$ (\autoref{eq:analytic_force_x}) --- a $4\times10^{-5}$
relative loss at $\delta=0.5^\circ$, negligible as a performance effect, while the $T\sin\delta$
that goes with it is not at all negligible as a \emph{moment}
(\autoref{sec:analysis_misaligned}). A \emph{commanded} gimbal, driven by a control law and the
only pitch authority a rocket has once it is off the rail, is left for later
(\autoref{sec:backlog}); $\delta$ here is simply an input to be swept, to see how sensitive the
airframe is to a manufacturing-grade misalignment before any control authority is added.

Propellant mass flow and the resulting CoG/moment-of-inertia travel are modelled in
\autoref{sec:mass}, and are carried by \autoref{sec:eom}: $m$, $l_B$ and $I$ are all functions of
time there.

\subsection{Mass, centre of gravity and moment of inertia}\label{sec:mass}
As propellant burns, mass leaves the rocket through the nozzle and both the centre of gravity and
the moment of inertia change over the burn. A minimal model treats the rocket as two lumped
masses: dry structure (mass $m_d$, fixed centroid at $l_d$ from the nose tip $A$) and remaining
propellant (mass $m_p(t)=m(t)-m_d$, centroid at $l_p$).

\subsubsection{Mass}
\begin{equation}
    \dot{m}(t) = -\frac{T_{ref}(t)}{I_{sp}\,g_0}
\end{equation}
or a measured mass-flow table, if available; $I_{sp}$ is the motor's specific impulse.

\subsubsection{Centre of gravity}
\begin{equation}
    l_B(t) = \frac{m_d\,l_d + m_p(t)\,l_p}{m_d + m_p(t)},
\end{equation}
the CoG distance from the nose, replacing the constant $l_B$ of \autoref{sec:geometry}.

\subsubsection{Moment of inertia}
\begin{equation}
    I(t) = I_d + m_d\bigl(l_B(t)-l_d\bigr)^2 + I_p(t) + m_p(t)\bigl(l_B(t)-l_p\bigr)^2,
\end{equation}
by the parallel-axis theorem, with $I_d,I_p(t)$ the dry structure's and remaining propellant's own
moments of inertia about their respective centroids.

$I_p(t)$ and $l_p$ depend on the propellant grain geometry and on how it recedes as it burns, so
they deserve elaborating rather than just flagging. Solid-motor grains commonly come in a few
standard families: \textbf{cylindrical} grains (a solid or centrally bored cylinder, burning on
the outer surface, or --- the common BATES configuration --- on the inner bore and, if
uninhibited, the end faces too); \textbf{tubular/annular} grains (a thin-walled, case-bonded tube
burning on the inner bore only); \textbf{end-burning} grains (burning only on one flat face, which
gives a nearly constant burn area and hence a nearly constant thrust, but a centroid that moves
along the grain quite differently from a radially-burning one); and \textbf{star or finocyl}
grains (a star- or wagon-wheel-shaped central perforation, chosen specifically to tailor the
burn-area-versus-time curve, and so the thrust curve $T_{ref}(t)$ itself). Each family gives a
different relationship between burn area and web thickness $A_b(w)$; combined with a regression
rate $\dot{w}=r_b(p)$ (typically Saint-Robert's law, $r_b=ap^n$), the instantaneous propellant
volume, its centroid $l_p(t)$, and its own moment of inertia $I_p(t)$ about that centroid all
follow by direct integration over the receding grain shape.

That is exactly the information the placeholder above is standing in for: holding $l_p$ constant
and taking $I_p(t)\propto m_p(t)$ assumes a uniform density and a shape that merely shrinks in
overall size, self-similarly, as it burns --- which no real grain does. The true centroid and the
true inertia ratio both drift over the burn, and by how much depends on which family above is
used and on the specific grain's own dimensions: bore and outer diameter, grain length, and, for
star grains, the number and shape of the points. None of that is fixed anywhere in this document
yet, so the placeholder stays in place until an actual grain geometry is chosen; this section can
be revisited once one is.
\todo{change of mass or moment of inertia due to mechanisms (e.g.\ deployable surfaces) is not
covered above}

\subsection{Aerodynamics}\label{sec:aero}
\subsubsection{Force magnitudes}
\begin{equation}
    A_x = qSC_A(\alpha,M), \qquad A_z = qSC_N(\alpha,M),
\end{equation}
\begin{equation} \label{eq:aero_magnitudes}
    A_m = qSd_{ref}C_m(\alpha,M),
\end{equation}
with $q$ the dynamic pressure, $S$ the reference area and $d_{ref}$ the reference length (in
practice the body cross-section and its diameter) --- the same $S,d_{ref}$ that
\autoref{sec:analysis} uses and the same $d_{ref}$ that forms the Reynolds number of
\autoref{sec:reynolds_mach}. $C_A,C_N,C_m$ are the axial-force, normal-force and pitching-moment
coefficients. $l_1=l_D-l_B$ is the axial distance from the centre of gravity $B$ to the
aerodynamic reference point $D$ of \autoref{sec:geometry}, positive with $D$ aft, so that
$\vec{BD}=-l_1\hat{x}_r$ and the caliber static margin of \autoref{sec:analysis_stability} is
$\mathrm{SM}=l_1/d_{ref}$.

\subsubsection{Dynamic pressure}
\begin{equation}
    q = \tfrac12\rho V^2,
\end{equation}
with $\rho$ the local air density (\autoref{sec:density}) and $V$ the relative airspeed. Both $V$
and the angle of attack $\alpha$ that feeds the coefficients above are derived from the rocket's
velocity relative to the surrounding air in \autoref{sec:relwind}.

\subsubsection{Coefficients}
$C_A,C_N,C_m$ and $l_1$ are not derived in this document from the airframe's shape: they are read
from a database, populated once per airframe from wind-tunnel testing, CFD sweeps, or a
semi-empirical panel-method tool, and interpolated at run time. In practice this means a
rectangular table over $(M,\alpha)$: bracket and clamp $M$ at the grid edges (no extrapolation
beyond the table's Mach range), interpolate linearly in $\alpha$ within each bracketing row, then
blend the two rows linearly across the $M$ bracket. Lacking such a table, a constant/linear
fallback is used instead:
\begin{equation}
    C_A = C_{A0}, \qquad C_N = C_{N\alpha}\alpha, \qquad C_m = C_{m\alpha}\alpha, \qquad l_1 = l_1^{(0)}.
\end{equation}
Implementation note: the fallback $C_N,C_m$ are evaluated as $C_{N\alpha}\sin\alpha$,
$C_{m\alpha}\sin\alpha$ in code (identical to the formula above for small $\alpha$, within $0.4\%$
at $8.6^\circ$) rather than the raw linear form. This is a numerical-robustness fix, not a
different physical model: $\alpha=\operatorname{atan2}(w,u)$ is discontinuous at its own branch
cut ($u<0$, $w$ crossing zero), which an unperturbed flight coasting tail-first through $V=0$ at
apogee lands on exactly, and the unbounded linear form turns that into a solver-breaking jump in
the moment every evaluation. $\sin\alpha$ removes the discontinuity (and, at large $\alpha$, is
arguably more correct besides: it is exactly zero at $\alpha=\pi$, matching the zero net side
force a body of revolution has flying exactly tail-first by symmetry).

The resulting force is resolved along the body axes directly --- \emph{not} rotated into the
relative-wind frame $\hat{u},\hat{n}$ of \autoref{sec:relwind} --- since $C_A,C_N$ are themselves
body-axis (axial/normal) coefficients, the same convention the underlying wind-tunnel/CFD data is
extracted in, and their $\alpha$-dependence is already carried by the table.

What such a database has to cover: angle of attack $\alpha$ (nonlinear at high $\alpha$ once flow
separates), Mach number $M$ (compressibility, transonic/supersonic drag rise), Reynolds number
$Re$ (boundary-layer state, viscous/skin-friction contribution to $C_A$), and geometry (fin
planform and cant, body-fin interference, surface finish) --- introducing exactly the $Re,M$ of
\autoref{sec:reynolds_mach}.
\todo{expand once the CFD/data-extraction pipeline that populates this database is defined; also
depends on control-surface extension not covered above (airbrake, aileron, etc.); might want to
include pitch rate (\autoref{sec:analysis_stability}'s $c_{m,q}$) and, for 3D, other stuff}

\subsubsection{Linearization of coefficients}
Coefficients can be linearised in $\alpha$ for the static-stability analysis of
\autoref{sec:analysis_stability}:
\begin{equation}
    c_z = c_{z,0} + \alpha c_{z,\alpha},
\end{equation}
\begin{equation}
    c_m = c_{m,0} + \alpha c_{m,\alpha},
\end{equation}
where $c_z\equiv C_N$ and $c_m\equiv C_m$: the lowercase form is kept purely because that is the
historical notation of that section, not a different quantity.
\todo{what about $c_d$ in general? A closed analytic drag-polar form exists for lifting surfaces
(airfoils, wings), induced drag growing with the square of the lift coefficient,}
\begin{equation}
    c_x = c_{x,0} + \frac{c_z^2}{e\pi\lambda},
\end{equation}
with $e$ the Oswald efficiency factor and $\lambda$ the aspect ratio. That, however, is
specifically an airfoil/wing result; what the equivalent closed form is for a body of revolution
(the rocket's own fuselage, which is not a lifting surface) is not resolved here, which is exactly
why the table-driven $C_A$ of \autoref{eq:aero_magnitudes} is used instead of a formula like this
one.

\subsection{Atmosphere}\label{sec:atmosphere}
Goal: given altitude $h$ (and, for the interpolated model, horizontal position), provide the local
air density $\rho$, together with everything needed to form the Reynolds and Mach numbers of the
relative flow (\autoref{sec:relwind}), and the local speed of sound $a$ and dynamic viscosity
$\mu$.

\subsubsection{Model tiers}\label{sec:atmos_tiers}
Three tiers, increasing in fidelity, all exposing the same $(\rho, T, p, a, \mu)$ interface:
\begin{enumerate}
    \item \textbf{ISA} (International Standard Atmosphere): the standard piecewise-linear
    temperature-lapse model (troposphere + lower stratosphere), pressure from the hydrostatic
    equation, density from the ideal gas law, speed of sound $a=\sqrt{\gamma R_{specific}T}$,
    viscosity from Sutherland's law.
    \item \textbf{Real atmosphere} (single station): measured/forecast sounding data ($T(h)$,
    $p(h)$ at one geographic location and time), interpolated in altitude only.
    \item \textbf{Real atmosphere} (multi-station, spatially interpolated): the same kind of
    measured data from several geographic locations, interpolated in altitude, horizontal
    position, and optionally time, to estimate local conditions along the actual ground track.
\end{enumerate}

\paragraph{ISA.} Geopotential altitude from geometric altitude above mean sea level $h_{msl}$:
\begin{equation}
    H = \frac{R_{isa}h_{msl}}{R_{isa}+h_{msl}}, \qquad R_{isa} = 6356766\,\mathrm{m}.
\end{equation}
Seven layers, base geopotential altitude $H_b$, lapse rate $L$, base temperature/pressure
$T_b, p_b$:
\begin{center}
\begin{tabular}{l c c c c}
Layer & $H_b$ [m] & $L$ [K/m] & $T_b$ [K] & $p_b$ [Pa] \\
troposphere  & 0     & $-0.0065$ & 288.15 & 101325   \\
tropopause   & 11000 & $0$       & 216.65 & 22632.06 \\
stratosphere & 20000 & $+0.0010$ & 216.65 & 5474.889 \\
stratosphere & 32000 & $+0.0028$ & 228.65 & 868.0187 \\
stratopause  & 47000 & $0$       & 270.65 & 110.9063 \\
mesosphere   & 51000 & $-0.0028$ & 270.65 & 66.93887 \\
mesosphere   & 71000 & $-0.0020$ & 214.65 & 3.956420 \\
\end{tabular}
\end{center}
Within the layer containing $H$, temperature is linear in geopotential altitude:
\begin{equation}
    T = T_b + L(H-H_b).
\end{equation}
For a layer with non-zero lapse rate the pressure follows the polytropic relation,
\begin{equation}
    p = p_b(T/T_b)^{-g_{isa}/(LR_d)},
\end{equation}
and for an isothermal layer ($L=0$) the exponential form instead,
\begin{equation}
    p = p_b\exp\!\left(\frac{-g_{isa}(H-H_b)}{R_dT_b}\right).
\end{equation}
Density, speed of sound and viscosity then follow from the ideal gas law and Sutherland's law:
\begin{equation}
    \rho = \frac{p}{R_dT}, \qquad a = \sqrt{\gamma R_dT}, \qquad
    \mu = \mu_0\left(\frac{T}{T_0}\right)^{3/2}\frac{T_0+S_\mu}{T+S_\mu},
\end{equation}
with $R_d = 287.0528\,\mathrm{J/(kg\,K)}$, $\gamma = 1.4$, $g_{isa} = 9.80665\,\mathrm{m/s^2}$ (the
ISA's own fixed value, independent of \autoref{sec:gravity}), $\mu_0 = 1.716\times10^{-5}\,\mathrm{Pa\,s}$,
$T_0 = 273.15\,\mathrm{K}$, $S_\mu = 110.4\,\mathrm{K}$.

\paragraph{Sounding.} One station's measured $(h_{msl}, p, T, u, v)$ profile, temperature and wind
interpolated linearly, pressure log-linearly:
\begin{equation}
    T(h) = \mathrm{lerp}(h;\{h_i\},\{T_i\}), \quad
    p(h) = \exp\mathrm{lerp}(h;\{h_i\},\{\ln p_i\}), \quad
    u(h),v(h) = \mathrm{lerp}(h;\{h_i\},\{u_i,v_i\}).
\end{equation}
$\rho, a, \mu$ are derived from the interpolated $T, p$ exactly as ISA above --- never interpolated
directly, to stay thermodynamically consistent. Below the lowest level the profile continues with
the sounding's own local lapse rate; above the highest, an ISA-shaped profile re-anchored at
$(h_{top}, T_{top}, p_{top})$ keeps $T, p$ continuous at the seam. Wind is held at the edge value
outside the data range.

\paragraph{Field.} $N$ stations projected onto a tangent plane centred at the launch site
$(\phi_0,\lambda_0)$, together with the rocket's own ground position $(x\sin\chi, x\cos\chi)$
($x$ = downrange coordinate, $\chi$ = launch heading, degrees clockwise from North, the direction
of $+\hat{x}_0$):
\begin{equation}
    E = R_e\cos\phi_0\,(\lambda-\lambda_0), \qquad N = R_e\,(\phi-\phi_0).
\end{equation}
The $N$ profiles are combined, weighted by inverse squared distance, into one profile:
\begin{equation}
    w_i = \frac{1}{d_i^2+\varepsilon^2}, \qquad \hat{w}_i = \frac{w_i}{\sum_j w_j}, \qquad
    (T, \ln p, u, v) = \sum_i \hat{w}_i\,(T_i, \ln p_i, u_i, v_i),
\end{equation}
then interpolated in altitude exactly as the sounding tier. A single station reduces every
$\hat{w}_i$ to 1 and reproduces the sounding tier exactly.

\subsubsection{Density}\label{sec:density}
Water vapour is lighter than dry air --- $M_v=18.02$ against $M_d=28.96\,\mathrm{g/mol}$ --- so at
a given pressure and temperature, replacing some of the air with vapour makes the mixture
\emph{less} dense, not more. Damp air is light air. The effect is at the percent level, and it is
one of the few corrections in this document that can be taken straight from measured data, because
radiosondes report dewpoint alongside temperature and pressure.

\paragraph{Partial pressures.} Treat the mixture as two ideal gases sharing a volume, dry air at
partial pressure $p-e$ and vapour at partial pressure $e$:
\begin{equation}
    \rho = \frac{p-e}{R_dT} + \frac{e}{R_vT}
         = \frac{p}{R_dT}\left[1 - \frac{e}{p}\left(1 - \epsilon\right)\right], \qquad
    \epsilon = \frac{R_d}{R_v} = \frac{287.0528}{461.5} = 0.6220.
\end{equation}
The bracket is the whole correction, and it is always less than one. Folding it into an effective
temperature instead of an effective density is the standard bookkeeping --- the \emph{virtual
temperature}, the temperature dry air would need in order to have the moist mixture's density at
the same pressure:
\begin{equation} \label{eq:virtual_temperature}
    T_v = \frac{T}{1 - \dfrac{e}{p}\left(1-\epsilon\right)}, \qquad
    \rho = \frac{p}{R_dT_v},
\end{equation}
so that every formula keeps the dry-air gas constant $R_d$ and only the temperature fed to it
changes. $T_v \geq T$ always, with equality for perfectly dry air.

\paragraph{Vapour pressure from dewpoint.} The dewpoint $T_d$ is \emph{defined} as the temperature
to which the parcel must be cooled, at constant pressure and mixing ratio, to reach saturation ---
so the actual vapour pressure of the parcel equals the \emph{saturation} vapour pressure evaluated
at the dewpoint, $e=e_s(T_d)$. That identity is what makes the dewpoint column directly usable.
For $e_s$ the Magnus form is used, with the Alduchov--Eskridge coefficients, $T_d$ in
$^\circ\mathrm{C}$ and $e_s$ in hPa:
\begin{equation} \label{eq:magnus}
    e_s(T_d) = 6.1094\,\exp\!\left(\frac{17.625\,T_d}{T_d + 243.04}\right).
\end{equation}
This form is accurate to better than $0.4\%$ over $-40$ to $+50\,^\circ\mathrm{C}$, which covers
every dewpoint a sounding reports. The factor of 100 converting hPa to Pa before forming $e/p$ is
the one easy mistake here.

\paragraph{Scope: measured tiers only.} The correction applies to the \textbf{sounding} and
\textbf{field} tiers, and \emph{not} to ISA. That is not an omission: the International Standard
Atmosphere is defined as dry air of fixed composition, and its pressure column is computed on that
basis --- injecting moisture into it would give a profile that is neither the standard nor a
measurement, the same argument that keeps $g_{isa}=9.80665$ inside the ISA regardless of which
gravity model \autoref{sec:gravity} selects. A sounding with no dewpoint data falls back to
$\rho=p/(R_dT)$ level by level, so the correction is present exactly where there is data to
support it.

\paragraph{Magnitude, and what is deliberately left dry.} At a warm, humid surface
($T=30\,^\circ\mathrm{C}$, $T_d=26\,^\circ\mathrm{C}$, $p=101325\,\mathrm{Pa}$) the correction is
$-1.25\%$ in density; at a temperate surface ($T=14.2\,^\circ\mathrm{C}$, $T_d=10\,^\circ\mathrm{C}$)
$-0.48\%$; at a cold, dry one ($T_d=-33\,^\circ\mathrm{C}$) $-0.017\%$; near the tropopause it is
below $10^{-5}$ and vanishes on its own without needing a cutoff, because $e_s$ collapses with
temperature. So this is a boundary-layer effect in practice --- which is where dynamic pressure is
largest during the early ascent, so it is not negligible for that reason. The speed of sound and
the viscosity are deliberately \emph{not} corrected: a consistent moist $a$ needs the mixture's own
$\gamma$ as well as its $R$, the two effects partly cancel, and what remains is $\sim0.2\%$ in $a$,
far below the uncertainty of the Mach-dependent coefficients it would feed
(\autoref{sec:aero}). The hydrostatic interpolation of $p$ between levels likewise stays on dry
$T$; strictly it should use $T_v$, but it is a within-layer interpolation of measured endpoints,
not an integration over a deep column.

\subsubsection{Reynolds and Mach numbers}\label{sec:reynolds_mach}
\begin{equation}
    Re = \frac{\rho V d_{ref}}{\mu}, \qquad M = \frac{V}{a},
\end{equation}
with $V$ from \autoref{sec:relwind}, $d_{ref}$ the reference length (the body diameter --- the
same $d_{ref}$ that normalises the aerodynamic coefficients, \autoref{sec:aero}), $\mu$ the
dynamic viscosity, and $a$ the local speed of sound, both from the tier above.

\subsection{Wind}\label{sec:wind}
Goal: given time $t$ and position (altitude $h$, and horizontal position for the interpolated
tier), provide the wind velocity $\vec{v}_w$ in $\mathcal{F}_0$, used in
$\vec{v}_{rel} = \vec{v} - \vec{v}_w$ (\autoref{sec:relwind}). Same three-tier structure as
\autoref{sec:atmos_tiers}:
\begin{enumerate}
    \item \textbf{Simple analytic model}: e.g.\ constant wind plus an altitude-dependent shear
    profile, optionally with seeded gusts.
    \item \textbf{Real wind data} (single station): measured/forecast wind profile $\vec{v}_w(h)$
    at one location and time.
    \item \textbf{Real wind data} (multi-station, spatially interpolated): wind profiles from
    multiple locations, interpolated in altitude, horizontal position, and optionally time.
\end{enumerate}

\paragraph{Flight-plane projection.} The launch heading $\chi$ (\autoref{sec:atmosphere}) fixes the
2D flight plane. A horizontal wind given as ENU components $(u,v)$, or as meteorological
speed/direction $(W,\psi)$ with $u=-W\sin\psi,\ v=-W\cos\psi$ ($\psi$ = the direction the wind
blows \emph{from}), projects onto it as
\begin{equation}
    w_x = u\sin\chi + v\cos\chi = -W\cos(\psi-\chi), \qquad
    w_y = u\cos\chi - v\sin\chi = -W\sin(\psi-\chi),
\end{equation}
with $w_x$ the downrange (along-plane) component feeding $\vec{v}_w$ and $w_z=0$ (no vertical wind
modelled). $w_y$, the cross-plane component, is computed and discarded: crosswind weathercocking
needs the out-of-plane freedom only a full 6DOF model has.

\paragraph{Analytic.} Power-law shear, constant direction, and optional deterministic gusts:
\begin{equation}
    W(h) = W_{ref}\left(\max(h,h_{ref})/h_{ref}\right)^{n}, \qquad \psi = \psi_0, \qquad
    \Delta W(t) = \sigma\sum_i w_i\sin(\omega_it+\varphi_i).
\end{equation}
$h$ is altitude above the local ground, $n\approx1/7$ the shear exponent (the classic open-terrain
value; smoother terrain is lower, rougher/urban/forested terrain higher, up to $\sim0.4$), valid
within the atmospheric surface/boundary layer only --- roughly the lowest few hundred metres to
1--2\,km depending on stability, above which the sounding tiers below should be used instead.
$\Delta W(t)$ is added to $W(h)$ before the projection above, then clamped non-negative, since $W$
is a speed. The $\{\omega_i,\varphi_i,w_i\}$ are drawn once and frozen --- the adaptive solver
samples the right-hand side at arbitrary, repeated and rejected times, so a fresh random draw per
call would make $\Delta W$ a non-function of $t$. $\sigma=0$ disables gusts entirely.

\paragraph{Sounding / field.} The same measured $(u,v)$ profiles and interpolation as
\autoref{sec:atmosphere}, projected through the formula above; never interpolated as $(W,\psi)$,
since direction wraps at $360^\circ$.

\paragraph{Vertical wind.} All tiers above return horizontal wind only ($w_z=0$);
updrafts/downdrafts are not modelled. Two standard families exist in the literature for deriving a
vertical component: (1) terrain-slope / mass-continuity orographic lift, $w\approx\vec{v}_h\cdot
\nabla h_{terrain}$, which needs a terrain elevation model this project does not have; (2)
thermal/convective updraft parameterisations (as used in the soaring/glider literature), which need
surface heat-flux or insolation data outside the radiosonde soundings used here. Neither fits the
current data pipeline (station soundings + a simple analytic profile) without adding a new data
source. This is deliberately left unimplemented for now, pending a look at which of the two
families (or a combination) is actually worth the added data dependency; revisit once that
research is done.

\subsection{Relative wind and angle of attack}\label{sec:relwind}
Let $\vec{v}$ be the rocket's velocity relative to $\mathcal{F}_0$, and $\vec{v}_w$ the wind
velocity relative to $\mathcal{F}_0$ (a model output, \autoref{sec:wind}). The rocket's velocity
relative to the surrounding air is
\begin{equation}
    \vec{v}_{rel} = \vec{v} - \vec{v}_w = u\hat{x}_r + w\hat{z}_r,
\end{equation}
with $u$ the axial (along-nose) component and $w$ the normal component, both expressed in
$\mathcal{F}_r$. The relative airspeed and angle of attack are
\begin{equation}
    V = \sqrt{u^2+w^2}, \qquad \alpha = \operatorname{atan2}(w, u),
\end{equation}
so that $u=V\cos\alpha$, $w=V\sin\alpha$. At $\alpha=0$ (nose-first into the relative wind)
$\vec{v}_{rel}=V\hat{x}_r$.

Define the unit vector along the relative wind and its in-plane normal (a $+90^\circ$ rotation of
$\hat{u}$ within $\mathcal{F}_r$):
\begin{equation}
    \hat{u} = \frac{\vec{v}_{rel}}{V} = \cos\alpha\,\hat{x}_r + \sin\alpha\,\hat{z}_r, \qquad
    \hat{n} = -\sin\alpha\,\hat{x}_r + \cos\alpha\,\hat{z}_r.
\end{equation}
At $\alpha=0$: $\hat{u}=\hat{x}_r$, $\hat{n}=\hat{z}_r$. $\hat{u},\hat{n}$ feed the dynamic pressure
$q$ and, through $\alpha$, the coefficient lookup in \autoref{sec:aero}; the aerodynamic force
itself is resolved along the fixed body axes $\hat{x}_r,\hat{z}_r$, \emph{not} along
$\hat{u},\hat{n}$ --- see \autoref{sec:aero}.

\subsection{Non-inertial frame corrections (Earth's rotation)}\label{sec:noninertial}
The WGS84 gravity model above already accounts for the Earth's rotation in the magnitude of $g$
(the static, equilibrium-shape effect, folded into $\gamma_0(\phi)$). A rotating Earth also
introduces two velocity/position-dependent fictitious forces into the translational equations of
motion, once $\mathcal{F}_0$ is reinterpreted as Earth-fixed rather than inertial:
\begin{equation}
    \vec{a}_{Coriolis} = -2\vec{\Omega}\times\vec{v}, \qquad
    \vec{a}_{centrifugal} = -\vec{\Omega}\times(\vec{\Omega}\times\vec{r}),
\end{equation}
with $\vec{\Omega}$ the Earth's angular velocity ($|\vec{\Omega}|=7.2921\times10^{-5}\,
\mathrm{rad/s}$) and $\vec{r},\vec{v}$ the rocket's position/velocity relative to the Earth-fixed
frame. The centrifugal term is a small, latitude-dependent correction, already absorbed into the
WGS84 $g(\phi,h)$ above when the launch site is the reference point; the Coriolis term is
genuinely new here and depends on the flight's own velocity and heading, not just the launch site.

\paragraph{Ground-frame basis.} Work in the local East--North--Up triad $(\hat{E},\hat{N},\hat{U})$
at the launch site $(\phi_0,\lambda_0)$, with $\phi_0$ the site's geodetic latitude (the same
quantity written $\phi$ in \autoref{sec:gravity}). The ground frame $\mathcal{F}_0$ is fixed by the
launch heading $\chi$ (degrees clockwise from North, the direction of $+\hat{x}_0$,
\autoref{sec:atmosphere} > Field) and by $\hat{z}_0$ pointing \emph{down} (\autoref{sec:frames}):
\begin{equation}
    \hat{x}_0 = \sin\chi\,\hat{E} + \cos\chi\,\hat{N}, \qquad
    \hat{y}_0 = \cos\chi\,\hat{E} - \sin\chi\,\hat{N}, \qquad
    \hat{z}_0 = -\hat{U}.
\end{equation}
This whole subsection is written in the \emph{fixed} ground frame $\mathcal{F}_0$, not in the
(possibly rotating) gravity frame $\mathcal{F}_g$: the $\vec{\Omega}$ decomposition below is
evaluated once, at the launch site, and does \emph{not} follow a rotating $\hat{z}_g(\varphi)$ ---
see the closing paragraph for what that costs. Nothing here is new: $\hat{x}_0$ is the downrange
direction used by the ground track $(x\sin\chi,x\cos\chi)$ of \autoref{sec:atmosphere} > Field,
and $\hat{y}_0$ is exactly the cross-plane direction whose wind component $w_y=u\cos\chi-v\sin\chi$
is computed and discarded in \autoref{sec:wind} > Flight-plane projection. The triad is
right-handed, $\hat{x}_0\times\hat{y}_0=\hat{z}_0$, which can be checked directly from the three
lines above ($(\sin\chi,\cos\chi,0)\times(\cos\chi,-\sin\chi,0)=(0,0,-1)$ in ENU components).
$\hat{y}_0$ points to the \emph{right} of the downrange direction.

\paragraph{Components of $\vec{\Omega}$.} The Earth's angular velocity points along the rotation
axis, i.e.\ toward the celestial north pole. At latitude $\phi_0$ that direction lies in the local
meridian plane (no East component), tilted from the local vertical by the co-latitude:
\begin{equation}
    \vec{\Omega} = \Omega\left(\cos\phi_0\,\hat{N} + \sin\phi_0\,\hat{U}\right), \qquad
    \Omega = 7.2921\times10^{-5}\,\mathrm{rad/s}.
\end{equation}
Projecting onto the ground-frame triad by taking dot products with the three basis vectors above:
\begin{equation}
    \Omega_x = \vec{\Omega}\cdot\hat{x}_0 = \Omega\cos\phi_0\cos\chi, \qquad
    \Omega_y = \vec{\Omega}\cdot\hat{y}_0 = -\Omega\cos\phi_0\sin\chi, \qquad
    \Omega_z = \vec{\Omega}\cdot\hat{z}_0 = -\Omega\sin\phi_0,
\end{equation}
so that
\begin{equation} \label{eq:omega_components}
    \vec{\Omega} = \Omega\left(\cos\phi_0\cos\chi\,\hat{x}_0
        - \cos\phi_0\sin\chi\,\hat{y}_0 - \sin\phi_0\,\hat{z}_0\right),
\end{equation}
with $\Omega_x^2+\Omega_y^2+\Omega_z^2=\Omega^2$ as required. Three sanity checks: at the North
pole ($\phi_0=90^\circ$) it reduces to $-\Omega\hat{z}_0$, i.e.\ straight up, for every $\chi$; on
the equator launching due North ($\phi_0=0,\chi=0$) to $+\Omega\hat{x}_0$, i.e.\ horizontal and
downrange; on the equator launching due East ($\phi_0=0,\chi=90^\circ$) to $-\Omega\hat{y}_0$,
which is horizontal and to the left of downrange --- North again, as it must be. The sign of
$\Omega_z$ is purely a bookkeeping consequence of this project's $z$-down convention: the rotation
vector points \emph{up} out of the northern hemisphere, hence negative along $\hat{z}_0$.

\paragraph{Which component of $\vec{\Omega}$ actually matters in 2D.} With motion confined to the
flight plane, $\vec{v}=v_x\hat{x}_0+v_z\hat{z}_0$ (no $\hat{y}_0$ component). Expanding the cross
product in the right-handed triad,
\begin{equation}
    \vec{\Omega}\times\vec{v} =
        \Omega_yv_z\,\hat{x}_0
        + \left(\Omega_zv_x - \Omega_xv_z\right)\hat{y}_0
        - \Omega_yv_x\,\hat{z}_0,
\end{equation}
so
\begin{equation} \label{eq:coriolis_split}
    \vec{a}_{Coriolis} = -2\vec{\Omega}\times\vec{v} =
        \underbrace{-2\Omega_y\left(v_z\hat{x}_0 - v_x\hat{z}_0\right)}_{\text{in-plane}}
        \;\underbrace{-\,2\left(\Omega_zv_x - \Omega_xv_z\right)\hat{y}_0}_{\text{out-of-plane}}.
\end{equation}
This is worth reading carefully, because it inverts the intuitive move. The in-plane Coriolis
acceleration is produced \emph{entirely} by $\Omega_y$ --- the component of $\vec{\Omega}$
\emph{normal} to the flight plane --- while $\Omega_x$ and $\Omega_z$, the components lying
\emph{inside} the plane, produce a purely out-of-plane acceleration. ``Project $\vec{\Omega}$ onto
the flight plane and keep that part'' would therefore keep exactly the half that does nothing
in-plane and throw away the half that does everything.

\paragraph{Simplification (2D closure).} Only the in-plane bracket of \autoref{eq:coriolis_split}
is retained; the $\hat{y}_0$ term is dropped. This is \emph{not} a smallness argument --- the
discarded lateral acceleration is of the same order as the retained one, and for some geometries
larger. It is dropped because the model has no lateral degree of freedom to receive it: the
out-of-plane Coriolis deflection (the familiar right-hand deflection in the northern hemisphere) is
physically real, and a 2D planar model simply cannot represent it. This is the same closure already
applied to the crosswind component $w_y$ in \autoref{sec:wind}, and it is lifted by the same
future work (a full 6DOF model), not by anything in this section.

Two consequences follow immediately and should be stated so results are not over-read: $\Omega_y=
-\Omega\cos\phi_0\sin\chi$ vanishes for a launch due North or due South ($\chi=0^\circ,180^\circ$)
and at the poles, so the 2D model predicts \emph{zero} Coriolis effect there while the true
Coriolis effect is maximal and entirely lateral; and $|\Omega_y|$ is largest for an East/West
launch near the equator, where the effect the 2D model \emph{can} see is also the largest it ever
is.

\paragraph{In-plane Coriolis term.} Substituting $\Omega_y$ into the retained bracket,
\begin{equation} \label{eq:coriolis_2d}
    \vec{a}_{Coriolis}^{\,2D} = 2\Omega\cos\phi_0\sin\chi\left(v_z\,\hat{x}_0 - v_x\,\hat{z}_0\right),
\end{equation}
i.e.\ componentwise, with the state variables of \autoref{sec:eom} ($v_x=\dot{x}$ downrange,
$v_z=\dot{z}$ \emph{downward}, so an ascending rocket has $v_z<0$):
\begin{equation}
    a_{Coriolis,x} = +2\Omega\cos\phi_0\sin\chi\;\dot{z}, \qquad
    a_{Coriolis,z} = -2\Omega\cos\phi_0\sin\chi\;\dot{x}.
\end{equation}
The two classical checks both come out right. A rocket rising vertically at the equator with an
Eastward $+\hat{x}_0$ ($\phi_0=0$, $\chi=90^\circ$, $\dot{z}<0$) gets $a_x<0$: it drifts
\emph{West}, the standard counterpart of the eastward deflection of a falling body. A body moving
East ($\dot{x}>0$) gets $a_z<0$, i.e.\ \emph{upward} --- the E\"otv\"os effect, eastward motion
reducing apparent weight.

\paragraph{Residual centrifugal term.} Write the rocket's geocentric position as
$\vec{r}=\vec{R}_0+\vec{s}$, with $\vec{R}_0$ from the Earth's centre to the launch site and
$\vec{s}=x\hat{x}_0+z\hat{z}_0$ the simulation's own position state. Since
$-\vec{\Omega}\times(\vec{\Omega}\times\vec{a})=\Omega^2\vec{a}-(\vec{\Omega}\cdot\vec{a})
\vec{\Omega}$ is linear in $\vec{a}$, the centrifugal acceleration splits cleanly into a constant
part, already carried by the WGS84 Somigliana formula, and a residual:
\begin{equation}
    \vec{a}_{centrifugal} = \underbrace{\Omega^2\vec{R}_0-(\vec{\Omega}\cdot\vec{R}_0)\vec{\Omega}}
        _{\text{constant, already in }\gamma_0(\phi)}
        + \underbrace{\Omega^2\vec{s}-(\vec{\Omega}\cdot\vec{s})\vec{\Omega}}_{\text{residual}}.
\end{equation}
This residual is bounded by $\Omega^2|\vec{s}|=5.3\times10^{-9}\,|\vec{s}|$, i.e.\
$5\times10^{-5}\,\mathrm{m/s^2}$ even at $|\vec{s}|=10\,\mathrm{km}$ --- roughly three orders of
magnitude below the Coriolis term, and its \emph{vertical} part is in any case already folded into
the WGS84 free-air correction through the constant $m$ in \autoref{sec:gravity} (double-counting it
would be a small error in the wrong direction). Recommendation: implement the Coriolis term only,
and leave the residual centrifugal term documented here but unimplemented.

\paragraph{Moment.} Coriolis and centrifugal accelerations are body forces distributed in exact
proportion to mass, so their resultant acts at the centre of gravity $B$ and contributes
\emph{zero} moment about $B$. The one surviving coupling comes from the body's own pitch rate
($\vec{v}_i=\vec{v}_B+\vec{\omega}\times\vec{r}_i$ inside the rigid body), a gyroscopic torque of
order $2I\Omega\omega$, i.e.\ an angular acceleration of order $2\Omega\omega\approx
1.5\times10^{-4}\,\mathrm{rad/s^2}$ at $\omega=1\,\mathrm{rad/s}$. This is negligible against the
aerodynamic pitching moment and is dropped; the rotational equation of \autoref{sec:eom} is
unchanged.

\paragraph{Result to add to the equations of motion.} The translational equation of
\autoref{sec:eom} gains one term:
\begin{equation}
    m\vec{\ddot{x}} = \vec{W} + \vec{T} + \vec{A}_f + m\,\vec{a}_{Coriolis}^{\,2D},
\end{equation}
with $\vec{a}_{Coriolis}^{\,2D}$ from \autoref{eq:coriolis_2d}, evaluated directly in
$\mathcal{F}_0$ (no body-axis rotation needed --- it is already a ground-frame expression, like
$\vec{W}$). $\phi_0$ and $\chi$ are the launch-site latitude and heading already used by
\autoref{sec:gravity} and \autoref{sec:wind} respectively; no new parameter is introduced.

\paragraph{Magnitudes, and an honest consistency caveat.} The retained term scales as
$2\Omega\cos\phi_0\sin\chi\,V$: at $V=300\,\mathrm{m/s}$ it is at most
$4.4\times10^{-2}\,\mathrm{m/s^2}$ (equatorial East launch) and about
$2.8\times10^{-2}\,\mathrm{m/s^2}$ at $\phi_0=50^\circ$ --- some $0.3\%$ of $g$, accumulating to
order $10\,\mathrm{m}$ of displacement over a $30\,\mathrm{s}$ powered ascent. That is worth
having for impact-point prediction and is consistent with the centimetre-to-metre estimate quoted
in \autoref{sec:assumptions}. It is roughly the same order as the effective
$g\,x/R_e\approx1.5\times10^{-2}\,\mathrm{m/s^2}$ (at $x=10\,\mathrm{km}$) that gravity acquires
by pointing along the rotated local vertical of \autoref{sec:gravity} > Direction, so the two are
worth having together rather than one without the other; they remain \emph{independent} ---one is
a static, position-dependent effect on the direction of a real force, the other a dynamic,
velocity-dependent fictitious force --- and must not be multiplied together.

\paragraph{One residual inconsistency, and why it is left standing.} The decomposition
\autoref{eq:omega_components} is evaluated once, at the launch site's latitude $\phi_0$ and
heading $\chi$, and is \emph{not} re-evaluated as the rocket travels downrange --- so
$\vec{\Omega}$ is resolved in the fixed ground frame $\mathcal{F}_0$ while gravity, if the
radial/elliptical direction model is enabled, is resolved in the rotated frame $\mathcal{F}_g$.
Deliberately so: building a single consistently rotating frame would replace both corrections with
a curvilinear formulation the 2D state cannot carry (it would need the metric terms of motion over
a sphere, not just a rotated basis). The cost of leaving them independent is bounded and small.
Along the ground track the latitude changes by $x\cos\chi/R_e$ and the heading by the meridian
convergence, so $\Omega_y$ drifts by $O(\Omega\,x/R_e)$, i.e.\ a relative $3\times10^{-6}$ of the
retained Coriolis acceleration at $x=10\,\mathrm{km}$ and $3\times10^{-4}$ at $100\,\mathrm{km}$
--- five orders of magnitude smaller than the gravity-direction term the same $x/R_e$ produces,
because the Coriolis term is itself $10^{-3}$ of gravity to begin with.


\newpage
\section{Backlog and future work}\label{sec:backlog}
The complement of \autoref{sec:assumptions}: that section records what the present model
simplifies, this one what is deliberately left undone. Each entry says what the item is and ---
more usefully --- what has to exist before it can be written down properly, so that a reader can
tell a missing derivation from a missing input.

\begin{itemize}
    \item \textbf{Commanded gimbal control.} $\delta$ is, for now, deliberately kept as a constant
    input (\autoref{sec:thrust}) --- the point of the current pass was to see how a fixed
    manufacturing-type misalignment behaves (\autoref{sec:analysis_misaligned}) before adding any
    control authority. A \emph{commanded}, time-varying $\delta(t)$ driven by a control law is the
    next step once that is wanted. \emph{Waiting on:} a stated guidance objective (pitch
    programme? gravity turn? impact point?), actuator rate and deflection limits, and a sensor
    model to close the loop on.
    \item \textbf{Wing-body aerodynamics: analytic CoP, airframe, air brake.} Replace the
    tabulated or constant coefficients of \autoref{sec:aero} with a wing-body (Barrowman-style)
    build-up. Conceptually: the nose, the body, and each fin each contribute their own normal-force
    slope $C_{N\alpha,i}$ and their own centre-of-pressure location $X_i$ (a slender
    nose/body contributes a small, often near-zero net slope away from the nose in the classical
    linear treatment, aside from any boat-tail correction; each fin's contribution uses a
    corrected 2D lift-curve slope, an interference factor between fin and body, and an
    effective aspect ratio from its planform), and the assembly's centre of pressure is the
    $C_{N\alpha}$-weighted average of the component locations,
    $X_{cp}=\sum_iC_{N\alpha,i}X_i/\sum_iC_{N\alpha,i}$ --- so that $l_1(\alpha,M)$ is
    \emph{computed} from the airframe rather than read from a table, and an air brake enters as a
    switchable change to $C_A$ and to the CoP. The original Barrowman formulas are subsonic
    linear-theory results and would need a Mach correction (Prandtl--Glauert or, in the linear
    supersonic regime, Busemann theory) to extend past low subsonic speed.
    \emph{Waiting on:} an actual airframe --- nose shape, body diameter/length and any boat-tail,
    and fin planform (root/tip chord, span, sweep, thickness, count) --- which this document does
    not yet fix anywhere.
    \item \textbf{Dynamic stability, quantitatively.} The framework is now in the document
    (\autoref{sec:analysis_dynamic}): a pitch-damping derivative $c_{m,q}$ linearises into a
    damped short-period oscillator with natural frequency $\omega_n$ and damping ratio $\zeta$.
    What is missing is a value or model for $c_{m,q}$ itself. \emph{Waiting on:} either a measured
    value (forced-oscillation wind-tunnel testing, or unsteady CFD) or a full distributed-force
    derivation from the airframe geometry of the previous item.
    \item \textbf{Propellant grain geometry behind $I_p(t)$ and $l_p$.} The physical mechanism and
    the standard grain families (cylindrical/BATES, tubular, end-burning, star/finocyl) are now
    elaborated in \autoref{sec:mass}; what remains is choosing one and fixing its actual
    dimensions (bore/outer diameter, length, and, for star grains, the point geometry), from which
    $A_b(w)$, $l_p(t)$ and $I_p(t)$ follow by direct integration. Picking this up is the next
    conversation, not something to guess at here.
    \item \textbf{Vertical wind.} Two candidate families exist (terrain-slope/mass-continuity
    orographic lift, or a convective/thermal parameterisation, \autoref{sec:wind} > Vertical
    wind), each needing a data source (terrain elevation, or surface heat-flux/insolation data)
    this project does not currently have. Deliberately on hold pending research into which family,
    if either, is worth the added dependency --- not to be implemented before that research is
    done.
    \item \textbf{CFD/data-extraction pipeline for $C_A,C_N,C_m$.} \autoref{sec:aero} now states
    plainly that these coefficients, and $l_1$, are sourced from a database (wind-tunnel testing,
    CFD sweeps, or a semi-empirical panel-method tool) rather than derived in this document. What
    remains open is the pipeline that actually populates that database for a specific airframe ---
    the CFD run matrix or wind-tunnel campaign itself --- which is outside the scope of this
    document.
    \item Launch rail constraint, and recovery/parachute --- \autoref{sec:assumptions}.
    \item Out-of-plane motion: crosswind weathercocking (\autoref{sec:wind}) and the lateral
    Coriolis deflection (\autoref{sec:noninertial}). Both are discarded by the same 2D closure and
    both are lifted by the same future work, a full 6DOF model --- not by anything that can be
    done inside this document.
\end{itemize}

\end{document}
